\chapter{Zwartość lokalna}
\label{ch:topog-zwartosc-lokalna}

\noindent\textbf{Cel rozdziału.}
Zwartość lokalna pozwala w otoczeniu każdego punktu
pracować ze zbiorami zwartymi, chociaż cała przestrzeń
nie musi być zwarta. Zaczniemy od definicji i podstawowego
lematu o otoczeniach mających zwarte domknięcia,
a następnie udowodnimy lokalnie zwartą wersję
twierdzenia Baire'a. Wersja metryczna tego twierdzenia
oraz pojęcie przestrzeni Baire'a zostały omówione
w rozdziale~11.

\section{Zwartość lokalna i twierdzenie Baire'a}
\label{sec:topog-baire-lokalnie-zwarte}

\begin{definicja}[Zwartość lokalna]
\label{def:topog-baire-lokalna-zwartosc}
Przestrzeń topologiczna $X$ jest
\emph{lokalnie zwarta}, jeżeli każdy
$x\in X$ ma zwarte otoczenie:
istnieje zwarty $K\subseteq X$
taki, że $x\in\operatorname{int}_X K$.
\end{definicja}

\begin{lemat}[Otoczenia o zwartym domknięciu]
\label{lem:topog-baire-zwarte-domkniecie}
Niech $X$ będzie lokalnie zwartą przestrzenią Hausdorffa.
Dla każdego $x\in X$ i każdego otwartego otoczenia
$V$ punktu $x$ istnieje otwarty $W\ni x$
o zwartym domknięciu $\overline W^{\,X}\subseteq V$.
\end{lemat}
\begin{proof}
Z lokalnej zwartości wybieramy zwarty
$K\subseteq X$ z $x\in\operatorname{int}_X K$.
Przestrzeń $K$ jest zwarta i Hausdorffa,
więc normalna (rozdział~07). Zbiór
$O=V\cap\operatorname{int}_X K$ jest
otwartym otoczeniem $x$ w $K$.
Kryterium zwężania otoczeń w przestrzeni
regularnej (rozdział~05) daje zbiór $H$
otwarty w $K$ spełniający
$x\in H\subseteq\overline H^{\,K}\subseteq O$.
Ponieważ $H\subseteq\operatorname{int}_X K$,
jest on także otwarty w $X$: jeśli
$H=T\cap K$ dla otwartego $T\subseteq X$,
to $H=T\cap\operatorname{int}_X K$.
Przyjmujemy $W=H$. Zwartość $K$
i własność Hausdorffa zapewniają domkniętość
$K$ w $X$ (rozdział~07), toteż
$\overline W^{\,X}=\overline H^{\,K}\subseteq V$.
Domknięty podzbiór zwartego $K$
jest zwarty, co kończy dowód.
\end{proof}

\begin{twierdzenie}[Twierdzenie Baire'a --- wersja lokalnie zwarta]
\label{thm:topog-baire-lokalnie-zwarte}
Każda lokalnie zwarta przestrzeń Hausdorffa
jest przestrzenią Baire'a.
\end{twierdzenie}
\begin{proof}
Niech $U$ będzie niepustym zbiorem otwartym,
a $G_n$ --- otwartymi zbiorami gęstymi.
Wybieramy $x_1\in U\cap G_1$.
Lemat~\ref{lem:topog-baire-zwarte-domkniecie}
daje niepusty otwarty $W_1$ o niepustym
zwartym domknięciu $K_1\subseteq U\cap G_1$.

Indukcyjnie, mając niepusty otwarty $W_n$
o zwartym domknięciu $K_n$, wybieramy
$x_{n+1}\in W_n\cap G_{n+1}$
i na mocy lematu otrzymujemy niepusty otwarty
$W_{n+1}$, którego zwarte domknięcie
$K_{n+1}$ zawiera się w $W_n\cap G_{n+1}$.
Zbiory $K_n$ są więc niepuste, zstępujące,
domknięte w $X$ (gdyż $X$ jest Hausdorffa)
i wszystkie zawierają się w zwartym $K_1$.

Gdyby $\bigcap_{n\ge1}K_n=\varnothing$,
to otwarte w $K_1$ zbiory $K_1\setminus K_n$
tworzyłyby pokrycie $K_1$. Ze zwartości
istniałoby podpokrycie skończone,
a ponieważ $(K_n)$ jest zstępujący,
oznaczałoby to $K_N=\varnothing$ dla
pewnego $N$. Jest to sprzeczność.
Zatem $\bigcap_nK_n$ jest niepuste
i zawiera się w $U\cap\bigcap_nG_n$.
\end{proof}

